\chapter{Peer-to-Peer}
\label{bab:p2p}

\section*{Tujuan Pembelajaran}

\begin{enumerate}
  \item Menjelaskan perbedaan susunan \textit{peer-to-peer} dari susunan
    client-server beserta letak pengetahuan pada masing-masing susunan.
  \item Membedakan jaringan \textit{unstructured} dari jaringan
    \textit{structured}, lalu menentukan kapan masing-masing layak dipilih.
  \item Mengukur \textit{lookup cost}, \textit{failure resilience}, dan
    \textit{message scalability} pada ketiga susunan.
\end{enumerate}

\section{Pendahuluan}

\textit{Peer-to-peer} menata sebuah sistem menjadi sekumpulan \textit{node} yang
berkedudukan sama. Setiap \textit{node} melayani permintaan dan sekaligus mengajukan
permintaan, sehingga tidak ada pembagian tetap antara pihak yang melayani dan
pihak yang dilayani. Bentuk ini berlawanan dengan client-server pada Bab~1,
yaitu ketika satu pihak selalu melayani dan pihak lain selalu meminta.

Perbedaan kedudukan tersebut menimbulkan satu persoalan yang tidak dimiliki
client-server. Pada client-server, letak sebuah data diketahui sejak awal,
yaitu di server. Pada \textit{peer-to-peer}, tidak ada satu \textit{node} pun yang
mengetahui seluruh isi jaringan, sehingga menemukan sebuah data menjadi
pekerjaan penting. Fielding menempatkan keduanya sebagai gaya yang berbeda,
dan mencatat bahwa \textit{peer-to-peer} menukar kesederhanaan dengan
ketahanan~\cite{fielding2000architectural}.

Bab ini memakai domain yang sama dengan Bab~2 dan Bab~4, yaitu \textit{lookup} nilai
tukar mata uang. Yang berubah hanya satu hal, yaitu cara sebuah \textit{node}
menemukan kurs yang tidak dipegangnya sendiri. Ketiga cara yang dibandingkan
mengikuti kerangka yang disusun Richards dan
Ford~\cite{richards2020fundamentals}.

Empat istilah perlu disepakati lebih dahulu, sebab seluruh bab ini memakainya.
\textbf{Kunci} adalah nama sebuah data, sedangkan \textbf{nilai} adalah isinya.
Pada bab ini kuncinya pasangan mata uang, misalnya \texttt{USD/IDR}, dan
nilainya kurs pasangan itu, yaitu 17.753,00 rupiah. Tabelnya memuat 26 kunci,
diambil apa adanya dari Kurs Transaksi Bank Indonesia tanggal 18 September 2026.
Pekerjaan yang dibandingkan bab ini hanya satu, yaitu menemukan nilai dari
sebuah kunci, dan pekerjaan itu disebut \textit{lookup}.

Tiga istilah berikutnya menyangkut jalannya pertanyaan. \textit{Neighbor}
adalah \textit{node} yang dikenal langsung oleh sebuah \textit{node}, sehingga
pertanyaan dapat dikirim kepadanya tanpa perantara. \textit{Hop} adalah satu
kali perpindahan pertanyaan dari satu \textit{node} ke \textit{node}
berikutnya, dan jumlah \textit{hop} menunjukkan berapa kali pertanyaan itu
dioper sebelum tiba pada pemegang kuncinya. Pesannya dihitung terpisah, yaitu
satu pesan untuk setiap perpindahan, baik perpindahan pertanyaan maupun
perpindahan jawaban. Fungsi \textbf{hash} adalah fungsi
yang mengubah sembarang teks menjadi satu angka tetap, sehingga nama yang sama
selalu menghasilkan angka yang sama.

\section{\textit{Node} yang Sekaligus Klien dan Server}

Tiga susunan dibandingkan sepanjang bab ini, dan ketiganya menjawab pertanyaan
yang sama dengan pengetahuan yang berbeda. \textbf{Client-server} menyimpan
seluruh pengetahuan pada satu \textit{node} indeks. \textbf{\textit{Unstructured}}
menyebarkan kunci ke sembarang \textit{node}, sehingga tidak ada yang tahu letaknya.
\textbf{\textit{Structured}} menghitung letak kunci dari kuncinya sendiri, sehingga
letaknya ditemukan tanpa bertanya ke mana-mana.

Gambar~\ref{fig:tiga-susunan} menunjukkan ketiganya berdampingan. \textit{Node} yang
diberi warna gelap adalah \textit{node} yang mengetahui letak seluruh kunci.

\begin{figure}[htbp]
  \centering
  \begin{tikzpicture}[
    simpul/.style={circle, draw, minimum size=0.42cm, inner sep=0pt,
      font=\tiny, fill=blue!8},
    utama/.style={simpul, fill=praditagreen!25},
    garis/.style={thin, gray!70},
    judul/.style={font=\scriptsize\bfseries, align=center}
  ]
    % Susunan pertama, seluruh pengetahuan pada satu simpul indeks.
    \begin{scope}[shift={(0,0)}]
      \node[utama] (i) at (0,0) {I};
      \foreach \sudut in {30,90,150,210,270,330} {
        \node[simpul] (k\sudut) at (\sudut:1.15cm) {};
        \draw[garis] (i) -- (k\sudut);
      }
      \node[judul] at (0,-1.9) {client-server\\
        {\tiny\mdseries satu indeks tahu semua}};
    \end{scope}

    % Susunan kedua, sambungan acak tanpa simpul yang tahu letak kunci.
    \begin{scope}[shift={(4.2,0)}]
      \node[simpul] (a) at (-0.95,0.95) {};
      \node[simpul] (b) at (0.35,1.15) {};
      \node[simpul] (c) at (1.15,0.15) {};
      \node[simpul] (d) at (0.55,-0.95) {};
      \node[simpul] (e) at (-0.75,-0.75) {};
      \node[simpul] (f) at (-0.15,0.05) {};
      \draw[garis] (a) -- (b);
      \draw[garis] (b) -- (c);
      \draw[garis] (c) -- (d);
      \draw[garis] (d) -- (e);
      \draw[garis] (e) -- (a);
      \draw[garis] (f) -- (b);
      \draw[garis] (f) -- (d);
      \draw[garis] (f) -- (a);
      \node[judul] at (0,-1.9) {unstructured\\
        {\tiny\mdseries tidak ada yang tahu}};
    \end{scope}

    % Susunan ketiga, letak kunci dihitung sehingga penelusurannya terarah.
    \begin{scope}[shift={(8.4,0)}]
      \foreach \nomor/\sudut in {0/90,1/45,2/0,3/315,4/270,5/225,6/180,7/135} {
        \node[simpul] (c\nomor) at (\sudut:1.15cm) {};
      }
      \draw[garis] (c0) -- (c1);
      \draw[garis] (c1) -- (c2);
      \draw[garis] (c2) -- (c3);
      \draw[garis] (c3) -- (c4);
      \draw[garis] (c4) -- (c5);
      \draw[garis] (c5) -- (c6);
      \draw[garis] (c6) -- (c7);
      \draw[garis] (c7) -- (c0);
      \draw[-{Stealth[length=1.6mm]}, thick, praditagreen] (c6) -- (c2);
      \draw[-{Stealth[length=1.6mm]}, thick, praditagreen] (c2) -- (c0);
      \node[judul] at (0,-1.9) {structured\\
        {\tiny\mdseries letaknya dihitung}};
    \end{scope}
  \end{tikzpicture}
  \caption{Tiga susunan yang dibandingkan sepanjang bab ini. Yang membedakan
    bukan bentuk gambarnya, melainkan siapa yang mengetahui letak sebuah kunci}
  \label{fig:tiga-susunan}
\end{figure}

Analogi yang mendekati adalah mencari satu buku pegangan di sebuah asrama. Cara
pertama, seorang penjaga memegang daftar siapa meminjam buku apa, lalu semua
penghuni bertanya kepadanya. Cara kedua, penghuni bertanya kepada teman
sekamar, lalu teman itu bertanya lagi kepada temannya sampai bukunya ketemu.
Cara ketiga, seluruh penghuni menyepakati aturan bahwa buku berawalan huruf
tertentu selalu dititipkan di kamar bernomor tertentu, sehingga tidak ada yang
perlu ditanya.

Pembagian tersebut menentukan letak setiap keputusan.
Tabel~\ref{tab:tanggung-p2p} merangkum pengetahuan dan larangan ketiga susunan.

\begin{table}[htbp]
  \centering
  \caption{Pengetahuan yang dipegang dan yang sengaja tidak dipegang setiap
    susunan}
  \label{tab:tanggung-p2p}
  \begin{tabularx}{\textwidth}{|l|X|X|}
    \hline
    \textbf{Susunan} & \textbf{Yang diketahui} & \textbf{Yang tidak diketahui} \\
    \hline
    Client-server & Indeks tahu letak seluruh kunci & \textit{Node} penanya tidak mengenal satu pun \textit{node} lain \\
    \hline
    \textit{Unstructured} & \textit{Node} hanya mengenal beberapa \textit{neighbor}-nya & Tidak ada satu \textit{node} pun yang tahu letak kunci \\
    \hline
    \textit{Structured} & Letak kunci dapat dihitung dari kuncinya & Tidak ada \textit{node} yang mengenal seluruh anggotanya \\
    \hline
  \end{tabularx}
\end{table}

Ketiga susunan tersebut menemukan kunci dengan cara yang berbeda, dan ketiganya
dibahas satu per satu pada seksi berikutnya. Angka ongkos yang disebut di bawah
ini adalah ringkasan pengukuran pada Seksi~\ref{sec:ongkos-p2p}, yaitu dua
ratus \textit{lookup} pada jaringan delapan \textit{node}. Satu pesan dihitung
untuk setiap perpindahan pertanyaan, ditambah satu pesan untuk setiap
perpindahan jawaban. Pada susunan client-server, seluruh
kunci dititipkan kepada satu \textit{node} indeks. Peminta bertanya sekali
kepada indeks itu lalu menerima jawabannya, sehingga ongkosnya tetap dua pesan
dan satu \textit{hop}, berapa pun jumlah \textit{node}-nya.

Pada susunan \textit{unstructured}, kunci disebar acak dan tidak ada yang
mencatat letaknya. Peminta bertanya kepada seluruh \textit{neighbor}-nya, lalu
\textit{neighbor} meneruskan pertanyaan yang sama kepada \textit{neighbor}
mereka sampai batas empat \textit{hop}. Ongkosnya 19 pesan untuk satu
\textit{lookup} pada delapan \textit{node}.

Pada susunan \textit{structured}, letak kunci dihitung dari hash namanya,
sehingga tidak ada yang perlu ditanya lebih dahulu. Setiap \textit{node}
menyimpan daftar pendek berisi \textit{node} yang pantas dihubungi untuk
beberapa jarak berbeda, dan pertanyaan dioper ke \textit{node} terjauh pada
daftar itu yang belum melewati sasarannya. Daftar tersebut bernama
\textit{finger table}, dan isinya dibahas pada Seksi~\ref{sec:ring-p2p}.
Ongkosnya 4 pesan pada jaringan yang sama.

\textit{Node}-nya sendiri sangat sederhana, seperti pada
Listing~\ref{lst:simpul-p2p}. Isinya hanya titipan kunci, daftar
\textit{neighbor}, dan satu penanda hidup.

\begin{lstlisting}[style=PythonStyle, caption={Satu \textit{node} jaringan,
  tersedia di \berkas{sources/chapter07/jaringan.py}}, label=lst:simpul-p2p]
class Node:
  """Satu node jaringan beserta titipan kunci dan daftar neighbor-nya."""

  def __init__(self, nomor: int) -> None:
    self.nomor = nomor
    self.simpanan = {}
    self.tetangga = []
    self.hidup = True

  def ambil(self, kunci: str):
    """Mengembalikan nilai kurs yang dititipkan, atau None bila tidak ada."""
    if not self.hidup:
      return None
    return self.simpanan.get(kunci)
\end{lstlisting}

\textit{Node}-nya berupa objek di dalam satu proses, bukan proses terpisah
yang berbicara lewat soket. Pilihan itu diambil sebab yang diukur bab ini adalah
jumlah pesan dan jumlah \textit{hop}, bukan waktu. Soket hanya akan menambahkan
ragam penjadwalan sistem operasi tanpa mengubah satu pun kesimpulannya. Sebagai
gantinya, setiap perpindahan pesan dihitung satu per satu.

\section{\textit{Centralized Index}}

Susunan pertama menitipkan seluruh kunci kepada satu \textit{node}, yaitu \textit{node}
bernomor nol. \textit{Node} lain tidak menyimpan apa pun dan tidak saling mengenal,
sehingga satu-satunya jalan menemukan kurs adalah bertanya kepada indeks.

Ongkosnya tetap, yaitu dua pesan dan satu \textit{hop}, berapa pun jumlah
\textit{node}-nya. Kedua pesan itu dapat disebut satu per satu. Pesan pertama
adalah pertanyaan dari peminta kepada indeks, dan pesan kedua adalah jawaban
indeks kepada peminta. \textit{Hop}-nya satu, sebab pertanyaan tersebut hanya
menyeberangi satu sambungan sebelum tiba pada \textit{node} pemegang kuncinya.
Jawaban yang pulang tidak menambah \textit{hop}, tetapi tetap dihitung sebagai
satu pesan.

Listing~\ref{lst:cari-indeks} memperlihatkan \textit{lookup}-nya. Panjangnya hanya
beberapa baris, dan justru itulah daya tariknya.

\begin{lstlisting}[style=PythonStyle, caption={\textit{Lookup} lewat
  \textit{centralized index}, tersedia di
  \berkas{sources/chapter07/indeks.py}}, label=lst:cari-indeks]
def cari(self, nomor_peminta: int, kunci: str) -> jaringan.LookupResult:
  """Bertanya satu kali kepada indeks, lalu menunggu satu jawaban."""
  peminta = self.simpul[nomor_peminta]
  if not peminta.hidup:
    return jaringan.LookupResult(kunci, None, 0, 0)
  nilai = peminta.ambil(kunci)
  if nilai is not None:
    return jaringan.LookupResult(kunci, nilai, 0, 0)
  # Satu pesan pertanyaan berangkat, lalu satu pesan jawaban kembali.
  indeks = self.simpul[NOMOR_INDEKS]
  nilai = indeks.ambil(kunci)
  return jaringan.LookupResult(kunci, nilai, 2, 1)
\end{lstlisting}

Susunan ini bukan rekaan untuk keperluan bab. Napster menyimpan daftar \textit{file}
pada server pusat, sedangkan berkasnya sendiri berpindah langsung antar
pemakai. BitTorrent menempuh jalan serupa lewat \textit{tracker}, yaitu layanan
Hypertext Transfer Protocol (HTTP) yang menjawab dengan daftar \textit{node} pemegang
\textit{file}~\cite{cohen2008bittorrent}. Tabel~\ref{tab:indeks-nyata} merangkum
ketiganya.

\begin{table}[htbp]
  \centering
  \caption{Letak indeks pada tiga sistem beserta yang tersisa ketika indeksnya
    mati}
  \label{tab:indeks-nyata}
  \begin{tabularx}{\textwidth}{|l|X|X|}
    \hline
    \textbf{Sistem} & \textbf{Letak indeks} & \textbf{Ketika indeks mati} \\
    \hline
    Napster & Server pusat menyimpan daftar \textit{file} seluruh pemakai & \textit{Lookup} berhenti sama sekali \\
    \hline
    BitTorrent & \textit{Tracker} menjawab dengan daftar \textit{node} pemegang \textit{file} & Unduhan yang berjalan bertahan, unduhan baru tidak \\
    \hline
    Contoh bab ini & \textit{Node} nol menyimpan seluruh 26 kunci kurs & Nol dari 200 \textit{lookup} terjawab \\
    \hline
  \end{tabularx}
\end{table}

BitTorrent kemudian menambah jalan kedua, yaitu Distributed Hash Table (DHT)
berdasarkan Kademlia~\cite{maymounkov2002kademlia}. Penambahan itu menjawab
persoalan yang sama, yaitu \textit{tracker} yang mati membuat \textit{file} tidak lagi
dapat ditemukan. Persoalan itulah yang menimbulkan pertanyaan berikutnya. Bila
indeksnya dihapus sama sekali, dengan cara apa sebuah kunci masih dapat
ditemukan.

\section{Menyebar Pertanyaan ke \textit{Neighbor}}

Jawaban paling sederhana adalah bertanya kepada siapa saja yang dikenal, lalu
meminta mereka bertanya lagi. Cara itu disebut penyebaran atau
\textit{flooding}, dan dipakai Gnutella. Agar pertanyaannya tidak berputar
tanpa henti, setiap pertanyaan dibekali batas \textit{hop} atau Time To Live (TTL),
yaitu angka yang berkurang satu pada setiap penerusan.

Gambar~\ref{fig:tetangga-acak} menunjukkan susunan \textit{neighbor} pada contoh bab
ini. Sambungannya dibuat acak, dan kunci pun dititipkan ke sembarang \textit{node}.

\begin{figure}[htbp]
  \centering
  \begin{tikzpicture}[
    simpul/.style={circle, draw, minimum size=0.62cm, inner sep=0pt,
      font=\scriptsize, fill=blue!8},
    garis/.style={thin, gray!80},
    ket/.style={font=\tiny, inner sep=1.5pt}
  ]
    \node[simpul, fill=praditagreen!25] (n3) at (-1.5,1.0) {3};
    \node[simpul] (n2) at (0.5,0.4) {2};
    \node[simpul] (n1) at (2.1,2.0) {1};
    \node[simpul, fill=green!10] (n5) at (1.4,3.4) {5};
    \node[simpul] (n0) at (4.0,1.0) {0};
    \node[simpul] (n7) at (1.8,-0.6) {7};
    \node[simpul] (n4) at (4.2,-1.0) {4};
    \node[simpul] (n6) at (2.7,-2.1) {6};

    \draw[garis] (n0) -- (n1);
    \draw[garis] (n0) -- (n4);
    \draw[garis] (n0) -- (n7);
    \draw[garis] (n1) -- (n2);
    \draw[garis] (n1) -- (n4);
    \draw[garis] (n1) -- (n5);
    \draw[garis] (n1) -- (n7);
    \draw[garis] (n2) -- (n3);
    \draw[garis] (n3) -- (n7);
    \draw[garis] (n4) -- (n6);
    \draw[garis] (n4) -- (n7);
    \draw[garis] (n6) -- (n7);

    \node[ket, below=1pt of n3] {peminta};
    \node[ket, right=2pt of n5] {pemegang EUR/IDR, satu tetangga};
  \end{tikzpicture}
  \caption{Delapan \textit{node} dengan dua belas sambungan acak. \textit{Node} 5 hanya punya
    satu \textit{neighbor}, dan kelemahan itu baru terasa pada seksi ketahanan}
  \label{fig:tetangga-acak}
\end{figure}

Listing~\ref{lst:sebar-p2p} memperlihatkan inti penyebarannya. Dua aturan
dipakai persis seperti pada Gnutella, yaitu pertanyaan tidak dikembalikan
kepada pengirimnya, dan pertanyaan berulang dibuang penerimanya.

\begin{lstlisting}[style=PythonStyle, caption={Inti penyebaran pertanyaan,
  tersedia di \berkas{sources/chapter07/tetangga.py}}, label=lst:sebar-p2p]
while antrean:
  nomor, pengirim, jarak = antrean.popleft()
  if jarak >= BATAS_LOMPATAN:
    continue
  for tujuan in self.simpul[nomor].tetangga:
    if tujuan == pengirim:
      continue
    pesan += 1
    if tujuan in terlihat or not self.simpul[tujuan].hidup:
      continue
    terlihat.add(tujuan)
    temuan = self.simpul[tujuan].ambil(kunci)
    if nilai is None and temuan is not None:
      nilai = temuan
      lompatan = jarak + 1
    antrean.append((tujuan, nomor, jarak + 1))
\end{lstlisting}

Sebelum akibatnya dibahas, satu pertanyaan perlu dijawab lebih dahulu, yaitu
untuk apa cara semahal ini dipakai. Jawabannya ada pada kemampuan yang tidak
dimiliki kedua susunan lain. \textit{Query flooding} dapat menjawab pertanyaan
yang hanya menyebut sebagian nama, misalnya seluruh kunci yang memuat USD, sebab
setiap \textit{node} memindai titipannya sendiri. Kemampuan itu tidak menuntut
siapa pun memegang daftar, sehingga jaringan tanpa pemilik tetap dapat dicari
isinya. Perbandingannya dengan kedua susunan lain menyusul pada
Seksi~\ref{sec:konteks-p2p}, sesudah susunan ketiga diperkenalkan.

Aturan penyebaran tersebut punya dua akibat praktis. Pertama, penyebarannya tidak
berhenti walaupun kuncinya sudah ketemu, sebab tidak ada satu \textit{node} pun
yang tahu bahwa \textit{lookup} itu sudah selesai. Kedua, ongkos \textit{lookup} karena itu hampir tidak
bergantung pada letak kuncinya, melainkan pada jumlah sambungan di dalam
jaringan.

Akibat kedua itulah yang menimbulkan persoalan berikutnya. Pertanyaannya,
apakah pertanyaan benar-benar harus disebar ke semua arah sekaligus.

\section{\textit{Identifier Ring} dan \textit{Finger Table}}
\label{sec:ring-p2p}

Chord menjawab bahwa penyebaran itu tidak perlu, asalkan letak kunci dapat
dihitung~\cite{stoica2001chord}. Nama \textit{node} dan nama kunci sama-sama dipetakan
menjadi satu titik pada \textit{ring} oleh fungsi hash yang sama. Sebuah kunci lalu
dititipkan kepada penerusnya, yaitu \textit{node} pertama yang titiknya tidak
mendahului titik kunci itu.

Satu hal perlu diluruskan sejak awal. Kunci dan \textit{node} adalah dua hal
yang berbeda, dan kunci bukan nomor sebuah \textit{node}. Yang disamakan hanya
ruang angkanya, yaitu nama \textit{node} dan nama kunci sama-sama diubah fungsi
hash yang sama menjadi satu titik di antara 0 sampai 65535. Titik kunci hampir
tidak pernah jatuh tepat pada titik sebuah \textit{node}, dan justru karena
itulah aturan penerus dibutuhkan. Pada contoh bab ini, kunci EUR/IDR jatuh pada
titik 63784, sedangkan \textit{node} 0 berada pada titik 5425, sehingga kunci
tersebut dititipkan kepada \textit{node} 0 sebagai penerus pertama sesudah titik
63784.

Pertanyaan yang lebih mendasar adalah mengapa titik pada \textit{ring} itu
dibutuhkan. Letak kunci sebenarnya dapat dihitung tanpa \textit{ring}, misalnya
dengan sisa bagi. Aturannya satu baris, yaitu pengenal kunci dibagi jumlah
\textit{node}, lalu sisa pembagiannya dipakai sebagai nomor \textit{node}
pemiliknya. Pembagi pada aturan itu adalah jumlah \textit{node} itu sendiri.
Cara ini juga tidak membutuhkan indeks terpusat, dan jauh lebih mudah ditulis.

Bedanya baru muncul ketika satu \textit{node} bergabung, sebab pembaginya ikut
berubah dari delapan menjadi sembilan. Contohnya kunci \texttt{USD/IDR} yang
pengenalnya 40796. Dibagi delapan, sisanya empat, sehingga kuncinya dititipkan
kepada \textit{node} 4. Dibagi sembilan, sisanya delapan, sehingga kunci yang
sama harus pindah ke \textit{node} 8. Hal serupa terjadi pada hampir seluruh
kunci lainnya. Pada \textit{ring}, yang berpindah hanya kunci di dalam busur
milik \textit{node} baru. Sifat itulah yang disebut \textit{consistent hashing}, dan Karger beserta
rekannya merumuskannya untuk menempatkan salinan halaman
web~\cite{karger1997consistent}.

Selisih kedua cara itu dapat dihitung, dan Tabel~\ref{tab:perpindahan} memuat
hasilnya pada 26 kunci kurs bab ini. Angkanya diambil dari keluaran
\texttt{uji\_perpindahan.py}.

\begin{table}[htbp]
  \centering
  \caption{Kunci yang berpindah pemilik ketika satu \textit{node} bergabung}
  \label{tab:perpindahan}
  \begin{tabularx}{\textwidth}{|X|l|l|}
    \hline
    \textbf{Yang diukur} & \textbf{Sisa bagi} & \textbf{\textit{Identifier ring}} \\
    \hline
    Kunci berpindah, 8 menjadi 9 \textit{node} & 26 dari 26 & 1 dari 26 \\
    \hline
    Kunci berpindah, 16 menjadi 17 \textit{node} & 24 dari 26 & 2 dari 26 \\
    \hline
  \end{tabularx}
\end{table}

Perpindahan kunci bukan urusan ringan, sebab setiap kunci yang berpindah berarti
satu pemindahan data antar \textit{node}. Pada jaringan yang \textit{node}-nya
sering datang dan pergi, cara sisa bagi membuat hampir seluruh data berpindah
setiap kali satu \textit{node} bergabung. \textit{Ring} menahan perpindahan itu
pada satu busur saja, dan itulah gunanya.

Ruang pengenal pada contoh bab ini enam belas bit, yaitu 65536 titik. Chord
sungguhan memakai seratus enam puluh bit agar dua \textit{node} tidak pernah menempati
titik yang sama. Ruang delapan bit sempat dicoba, lalu ditinggalkan, sebab pada
32 \textit{node} dua di antaranya jatuh pada titik yang sama.
Gambar~\ref{fig:cincin-p2p} menunjukkan letak kedelapan \textit{node} beserta jalur
\textit{lookup} kunci EUR/IDR yang jatuh pada titik 63784.

\begin{figure}[htbp]
  \centering
  \begin{tikzpicture}[
    titik/.style={circle, fill=praditagreen, minimum size=0.17cm,
      inner sep=0pt},
    nama/.style={font=\tiny, inner sep=1.5pt},
    panah/.style={-{Stealth[length=2mm]}, thick, praditagreen}
  ]
    \draw[gray!60, thick] (0,0) circle (2.6cm);
    \node[font=\tiny, above] at (0,2.68) {titik 0};

    % Sudutnya dihitung searah jarum jam dari puncak, yaitu 90 derajat
    % dikurangi pengenal dibagi 65536 lalu dikali 360 derajat.
    \node[titik] (s0) at (60.2:2.6cm) {};
    \node[titik] (s2) at (134.1:2.6cm) {};
    \node[titik] (s7) at (150.8:2.6cm) {};
    \node[titik] (s4) at (165.6:2.6cm) {};
    \node[titik] (s1) at (207.1:2.6cm) {};
    \node[titik] (s5) at (214.0:2.6cm) {};
    \node[titik] (s3) at (266.4:2.6cm) {};
    \node[titik] (s6) at (296.8:2.6cm) {};
    \node[circle, draw=red!70, thick, minimum size=0.2cm, inner sep=0pt]
      at (99.6:2.6cm) {};

    \node[nama, anchor=south east, text=red!70] at (99.6:2.72cm)
      {EUR/IDR (63784)};
    \node[nama, anchor=south west] at (60.2:2.72cm) {node 0 (5425)};
    \node[nama, anchor=south east] at (134.1:2.72cm) {node 2 (57505)};
    \node[nama, anchor=east] at (150.8:2.72cm) {node 7 (54466)};
    \node[nama, anchor=east] at (165.6:2.72cm) {node 4 (51768)};
    \node[nama, anchor=east, yshift=5pt] at (207.1:2.72cm) {node 1 (44208)};
    \node[nama, anchor=east, yshift=-5pt] at (214.0:2.72cm) {node 5 (42962)};
    \node[nama, anchor=north] at (266.4:2.76cm) {node 3 (33413)};
    \node[nama, anchor=north west] at (296.8:2.72cm) {node 6 (27881)};

    \draw[panah] (s3) -- (s4);
    \draw[panah] (s4) -- (s2);
    \draw[panah] (s2) -- (s0);
    \node[font=\tiny, align=center] at (0.1,-0.15)
      {tiga hop\\ searah jarum jam};
  \end{tikzpicture}
  \caption{\textit{Identifier ring} enam belas bit beserta jalur
    \textit{lookup} EUR/IDR dari \textit{node} 3. Angka di dalam kurung adalah
    titik \textit{ring}-nya, dan kunci pada titik 63784 dititipkan kepada
    penerusnya, yaitu \textit{node} 0}
  \label{fig:cincin-p2p}
\end{figure}

Perhitungan titiknya dapat dicoba langsung. Nama di-hash dengan SHA-1, lalu
hasilnya diambil sisa baginya terhadap 65536, sehingga jatuh pada salah satu
titik ruang enam belas bit. Nama kunci dan nama \textit{node} memakai fungsi
yang sama, sehingga \texttt{EUR/IDR} jatuh pada titik 63784 sedangkan
\texttt{simpul-0} pada titik 5425. Pemiliknya dicari searah jarum jam, yaitu
\textit{node} pertama yang titiknya tidak kurang dari 63784. Titik \textit{node}
terbesar pada contoh ini hanya 57505, sehingga pencariannya melewati titik 0 lalu
berhenti di \textit{node} 0.

Letak kunci memang dapat dihitung, tetapi tidak satu \textit{node} pun
menyimpan daftar seluruh \textit{node}. Yang pasti diketahui setiap \textit{node}
hanya penerusnya, yaitu \textit{node} terdekat searah jarum jam. Pertanyaannya
karena itu dioper dari satu \textit{node} ke \textit{node} berikutnya sampai tiba
pada pemegang kunci, dan setiap operan itulah yang dihitung sebagai satu
\textit{hop}.

Mengoper lewat penerus saja memang sampai, tetapi jalannya merayap. Pada delapan
\textit{node}, satu \textit{lookup} dapat menempuh tujuh \textit{hop}, yaitu satu
kurang daripada jumlah \textit{node}-nya. Pada sejuta \textit{node}, jalan
merayap seperti itu jelas tidak terpakai.

Jalan pintasnya disebut \textit{finger table}. Bentuknya buku alamat kecil di
dalam setiap \textit{node}, berisi enam belas baris, dan setiap baris menyimpan
satu nama \textit{node} yang pantas dihubungi untuk satu jarak tertentu. Baris
pertama memuat \textit{node} terdekat searah jarum jam, baris berikutnya untuk
jarak dua kali lipat, lalu empat kali lipat, dan seterusnya sampai separuh
\textit{ring}. Di dalamnya tidak ada daftar seluruh \textit{node}, sehingga
tabelnya tetap enam belas baris walaupun jaringannya berisi sejuta
\textit{node}.

Gunanya satu, yaitu menahan jumlah \textit{hop} agar tidak tumbuh mengikuti
jumlah \textit{node}.
\textit{Finger} ke-$i$ menunjuk penerus dari titik dirinya ditambah dua pangkat
$i$, sehingga jarak jalan pintasnya berlipat dua, yaitu 1, 2, 4, 8, dan
seterusnya sampai separuh \textit{ring}. Setiap \textit{hop} karena itu memangkas
sedikitnya separuh sisa jarak.

Selisihnya baru terasa pada jaringan besar. Lewat penerus saja, satu
\textit{lookup} rata-rata menempuh separuh jumlah \textit{node}, sehingga pada
sejuta \textit{node} angkanya sekitar lima ratus ribu \textit{hop}. Lewat
\textit{finger table}, angkanya mengikuti logaritma dua jumlah \textit{node},
sehingga sejuta \textit{node} cukup ditempuh sekitar dua puluh \textit{hop}.

Selisih kedua jalan itu dapat diukur, yaitu dengan mematikan sementara pemakaian
\textit{finger table} pada \berkas{sources/chapter07/cincin.py}.
Tabel~\ref{tab:jalan-pintas} memuat hasilnya pada dua ratus \textit{lookup} yang
sama.

\begin{table}[htbp]
  \centering
  \caption{Penelusuran \textit{ring} dengan dan tanpa \textit{finger table},
    dua ratus \textit{lookup} pada delapan \textit{node}}
  \label{tab:jalan-pintas}
  \begin{tabularx}{\textwidth}{|X|l|l|}
    \hline
    \textbf{Yang diukur} & \textbf{Penerus saja} & \textbf{\textit{Finger table}} \\
    \hline
    \textit{Hop} median & 4 & 2 \\
    \hline
    \textit{Hop} terbesar & 7 & 4 \\
    \hline
    Pesan median & 8 & 4 \\
    \hline
    \textit{Hop} EUR/IDR dari \textit{node} 3 & 6 & 3 \\
    \hline
  \end{tabularx}
\end{table}

Jalur pada Gambar~\ref{fig:cincin-p2p} karena itu hanya tiga \textit{hop}.
Pertama, \textit{node} 3 memakai \textit{finger} ke-14 lalu mendarat di
\textit{node} 4 pada titik 51768. Kedua, \textit{node} 4 memakai \textit{finger}
ke-12 lalu mendarat di \textit{node} 2 pada titik 57505. Ketiga, titik 63784
sudah berada di busur antara \textit{node} 2 dan penerusnya, sehingga
pertanyaannya diserahkan langsung kepada \textit{node} 0.

Tabel~\ref{tab:tabel-jari} memperlihatkan isi \textit{finger table}
\textit{node} 3 pada contoh bab ini.

\begin{table}[htbp]
  \centering
  \caption{\textit{Finger table} \textit{node} 3, yang titik cincinnya 33413}
  \label{tab:tabel-jari}
  \begin{tabularx}{\textwidth}{|l|l|l|X|}
    \hline
    \textbf{\textit{Finger} ke-$i$} & \textbf{Titik dicari} & \textbf{Penerusnya} & \textbf{Titik penerus} \\
    \hline
    0 sampai 13 & 33414 sampai 41605 & \textit{node} 5 & 42962 \\
    \hline
    14 & 49797 & \textit{node} 4 & 51768 \\
    \hline
    15 & 645 & \textit{node} 0 & 5425 \\
    \hline
  \end{tabularx}
\end{table}

Empat belas \textit{finger} pertama menunjuk \textit{node} yang sama, dan
kejanggalan itu memang wajar pada jaringan kecil. Jarak dua pangkat nol sampai
dua pangkat tiga belas seluruhnya belum melampaui \textit{node} 5, sehingga
penerusnya tidak berubah. Baris terakhir memperlihatkan tepi \textit{ring}-nya,
sebab titik 33413 ditambah 32768 melewati 65535 lalu jatuh kembali ke titik 645.
\textit{Finger table} baru terasa gunanya ketika jumlah \textit{node}-nya banyak,
dan Latihan~3 memeriksa hal itu dengan angka.

Listing~\ref{lst:telusuri-p2p} memperlihatkan penelusurannya.
\textit{Finger} terjauh yang belum melewati titik sasaran dipilih lebih dahulu,
sebab jalan pintas itulah yang memangkas sisa jarak paling banyak.

\begin{lstlisting}[style=PythonStyle, caption={Penelusuran \textit{ring}
  lewat \textit{finger table}, tersedia di
  \berkas{sources/chapter07/cincin.py}}, label=lst:telusuri-p2p]
sekarang = peminta
lompatan = 0
while lompatan < BATAS_LOMPATAN:
  penerus = self.penerus_hidup(sekarang)
  if penerus is None:
    break
  if di_antara(sasaran, sekarang.pengenal, penerus.pengenal):
    lompatan += 1
    nilai = penerus.ambil(kunci)
    return jaringan.LookupResult(kunci, nilai, lompatan * 2, lompatan)
  berikut = self.jari_terdekat(sekarang, sasaran) or penerus
  sekarang = berikut
  lompatan += 1
\end{lstlisting}

\section{Ongkos Satu \textit{Lookup}}
\label{sec:ongkos-p2p}

Sebelum angkanya dibandingkan, satu hal perlu dipastikan lebih dahulu.
Ketiga susunan dibangun dari satu kerangka yang sama, yaitu kelas \texttt{Network}
pada \berkas{sources/chapter07/jaringan.py}. Kerangka itu menyiapkan \textit{node},
menyediakan pengacak dengan \textit{seed} tetap, lalu menyerahkan dua metode
kepada turunannya. Metode pertama \texttt{susun}, yaitu cara menitipkan kunci.
Metode kedua \texttt{cari}, yaitu cara menemukan kunci itu kembali. Bahan
lainnya sama persis pada ketiga susunan, yaitu 26 kunci kurs, daftar peminta,
dan \textit{seed} pengacaknya. Karena yang berbeda hanya kedua metode tersebut,
selisih jumlah pesan dan jumlah \textit{hop} antara ketiganya benar-benar
berasal dari cara menitipkan dan cara mencari.

Tujuh kelas dipakai sepanjang bab ini, dan Tabel~\ref{tab:kelas-p2p}
menyebutkan peran masing-masing. Tiga kelas pertama menjadi kerangka bersama,
sedangkan empat berikutnya mewujudkan ketiga susunan yang dibandingkan.

\begin{table}[htbp]
  \centering
  \caption{Kelas yang dipakai ketiga susunan beserta perannya}
  \label{tab:kelas-p2p}
  \begin{tabularx}{\textwidth}{|l|X|}
    \hline
    \textbf{Kelas} & \textbf{Perannya} \\
    \hline
    \texttt{Network} & Kerangka bersama, menyiapkan \textit{node} lalu menyerahkan \texttt{susun} dan \texttt{cari} kepada turunannya \\
    \hline
    \texttt{Node} & Satu \textit{node} beserta titipan kunci, daftar \textit{neighbor}, dan penanda hidup \\
    \hline
    \texttt{LookupResult} & Catatan satu \textit{lookup}, yaitu nilai temuan, jumlah pesan, dan jumlah \textit{hop} \\
    \hline
    \texttt{IndexNetwork} & Susunan client-server, seluruh kunci dititipkan kepada satu \textit{node} indeks \\
    \hline
    \texttt{FloodingNetwork} & Susunan \textit{unstructured}, pertanyaan disebar ke seluruh \textit{neighbor} \\
    \hline
    \texttt{RingNetwork} & Susunan \textit{structured}, letak kunci dihitung pada \textit{identifier ring} \\
    \hline
    \texttt{RingNode} & Turunan \texttt{Node} yang membawa titik \textit{ring} beserta \textit{finger table}-nya \\
    \hline
  \end{tabularx}
\end{table}

Gambar~\ref{fig:kelas-p2p} menunjukkan hubungan ketujuh kelas tersebut.
Perhatikan bahwa \textit{finger table} hanya dimiliki \texttt{RingNode},
sedangkan dua susunan lain memakai \texttt{Node} apa adanya.

\begin{figure}[htbp]
  \centering
  \includegraphics[width=\textwidth, height=0.40\textheight,
    keepaspectratio]{kelas-p2p.pdf}
  \caption{Satu kerangka bersama beserta tiga susunan yang menurunkannya. Hanya
    \texttt{susun} dan \texttt{cari} yang ditimpa masing-masing susunan}
  \label{fig:kelas-p2p}
\end{figure}

Ketiga susunan kini dapat diadu dengan bahan yang sama. Dua ratus \textit{lookup}
dijalankan, dan daftar pemintanya sama persis pada ketiganya. Peminta tidak
pernah diambil dari \textit{node} nol, sebab pada susunan client-server \textit{node} itulah
indeksnya, dan indeks tidak bertanya kepada dirinya sendiri.

Gambar~\ref{fig:sekuens-unstructured} dan Gambar~\ref{fig:sekuens-structured}
menunjukkan satu \textit{lookup} yang sama lewat dua jalur berbeda. Kuncinya
EUR/IDR dan pemintanya \textit{node} 3 pada keduanya, sedangkan jumlah
\textit{hop}-nya sama, yaitu tiga.

\begin{figure}[htbp]
  \centering
  \includegraphics[width=\textwidth, height=0.50\textheight,
    keepaspectratio]{sekuens-p2p-unstructured.pdf}
  \caption{Jalur \textit{unstructured}. Pertanyaan disebar ke seluruh
    \textit{neighbor}, sehingga 20 pesan melintas untuk satu \textit{lookup}}
  \label{fig:sekuens-unstructured}
\end{figure}

\begin{figure}[htbp]
  \centering
  \includegraphics[width=0.86\textwidth, height=0.42\textheight,
    keepaspectratio]{sekuens-p2p-structured.pdf}
  \caption{Jalur \textit{structured}. Penelusurannya lewat \textit{finger
    table}, sehingga cukup 6 pesan dengan jumlah \textit{hop} yang sama}
  \label{fig:sekuens-structured}
\end{figure}

Hasil dua ratus \textit{lookup} tercantum pada Tabel~\ref{tab:ongkos-p2p}. Angkanya
diambil apa adanya dari keluaran \texttt{ukur\_pencarian.py}.

\begin{table}[htbp]
  \centering
  \caption{Ongkos \textit{lookup} pada delapan \textit{node}, dua ratus \textit{lookup}}
  \label{tab:ongkos-p2p}
  \begin{tabularx}{\textwidth}{|X|l|l|l|}
    \hline
    \textbf{Yang diukur} & \textbf{Client-server} & \textbf{\textit{Unstructured}} & \textbf{\textit{Structured}} \\
    \hline
    Pesan median & 2 & 19 & 4 \\
    \hline
    Pesan persentil ke-95 & 2 & 20 & 6 \\
    \hline
    Pesan terbesar & 2 & 20 & 8 \\
    \hline
    \textit{Hop} median & 1 & 2 & 2 \\
    \hline
    \textit{Hop} terbesar & 1 & 3 & 4 \\
    \hline
    Dijawab \textit{node} sendiri & 0 & 25 & 29 \\
    \hline
    Sisi jaringan & 7 & 12 & 13 \\
    \hline
  \end{tabularx}
\end{table}

Tiga hal layak dibaca dari tabel tersebut. Pertama, kedua susunan
\textit{peer-to-peer} memakai jumlah sambungan yang hampir sama, yaitu dua
belas melawan tiga belas, sehingga selisih pesannya benar-benar berasal dari
cara mencari. Kedua, susunan \textit{unstructured} memakai hampir lima kali pesan
susunan \textit{structured}, padahal \textit{hop} mediannya sama. Ketiga, sebaran susunan
\textit{unstructured} nyaris rata, yaitu median 19 dan terbesar 20, sebab ongkosnya
memang tidak bergantung pada letak kunci.

Baris kelima memuat kejanggalan yang perlu dijelaskan. Sebanyak 25 \textit{lookup}
pada susunan \textit{unstructured} dijawab tanpa satu pun pesan, sebab kuncinya
kebetulan dipegang \textit{node} peminta sendiri. Angka itu menarik median pesannya
turun sedikit, dan angka sejenis muncul pula pada susunan \textit{structured} sebanyak
29 kali. Susunan client-server tidak pernah mengalaminya, sebab seluruh
kuncinya memang berada di tempat lain.

\section{Ketika Satu \textit{Node} Mati}

Ongkos \textit{lookup} baru separuh cerita. Pertanyaan yang lebih menentukan adalah
berapa banyak \textit{lookup} yang masih terjawab ketika satu \textit{node} keluar dari
jaringan. Percobaannya diulang untuk setiap \textit{node}, yaitu setiap \textit{node}
bergiliran menjadi \textit{node} yang mati.

Kasus berikut memperlihatkan akibatnya. Sebuah tim membangun layanan kurs
internal memakai susunan \textit{centralized index}, dan layanannya dipakai delapan kantor
cabang. Tabel kursnya disimpan pada satu mesin di kantor pusat, sebab cara itu
paling mudah diawasi. Pada suatu Jumat sore, mesin tersebut dimatikan untuk
pemeliharaan yang dijadwalkan dua jam. Seluruh cabang seketika kehilangan kurs,
padahal masing-masing sudah membaca tabel yang sama berkali-kali. Tidak satu
cabang pun menyimpan salinannya, sebab susunannya memang tidak menuntut hal
itu, dan pemeliharaan dua jam berubah menjadi dua jam tanpa transaksi.

Tabel~\ref{tab:ketahanan-p2p} memuat hasil percobaan nyatanya. Angkanya diambil
dari keluaran \texttt{uji\_ketahanan.py}.

\begin{table}[htbp]
  \centering
  \caption{\textit{Lookup} yang masih terjawab ketika satu \textit{node} mati}
  \label{tab:ketahanan-p2p}
  \begin{tabularx}{\textwidth}{|X|l|l|l|}
    \hline
    \textbf{Yang diukur} & \textbf{Client-server} & \textbf{\textit{Unstructured}} & \textbf{\textit{Structured}} \\
    \hline
    Seluruh \textit{node} hidup & 200 dari 200 & 200 dari 200 & 200 dari 200 \\
    \hline
    \textit{Node} 0 mati & 0,0 persen & 94,0 persen & 86,0 persen \\
    \hline
    Rata-rata satu \textit{node} mati & 87,5 persen & 83,5 persen & 87,9 persen \\
    \hline
    Terburuk satu \textit{node} mati & 0,0 persen & 55,2 persen & 71,5 persen \\
    \hline
    Nomor \textit{node} terburuk & 0 & 1 & 6 \\
    \hline
    Kunci pada \textit{node} terburuk & 26 & 4 & 7 \\
    \hline
  \end{tabularx}
\end{table}

Baris ketiga tabel tersebut memuat kejanggalan yang paling penting pada bab
ini. Menurut rata-rata, susunan client-server mengungguli susunan
\textit{unstructured}, yaitu 87,5 persen melawan 83,5 persen. Kesimpulan itu keliru, dan
yang keliru bukan angkanya melainkan ukurannya. Tujuh dari delapan giliran pada
susunan client-server tidak menghilangkan apa pun, sebab tujuh \textit{node} itu
memang tidak menyimpan kunci. Giliran kedelapan menghilangkan segalanya, dan
rata-rata menyembunyikan giliran tersebut.

Baris keempat memperlihatkan hal yang tidak dapat disembunyikan. Susunan
client-server jatuh ke nol persen, sedangkan kedua susunan
\textit{peer-to-peer} tetap menjawab lebih dari separuh \textit{lookup}. Inilah yang
disebut titik kegagalan tunggal, dan inilah yang dibeli dengan ongkos pesan
yang lebih mahal pada seksi sebelumnya.

Kedua susunan \textit{peer-to-peer} pun tidak sama kuatnya, dan sebabnya
berbeda. Susunan \textit{unstructured} jatuh paling dalam ketika \textit{node} 1 mati,
sebab \textit{node} 5 hanya bertetangga kepadanya, sehingga enam kunci milik \textit{node} 5
ikut tidak terjangkau. Susunan \textit{structured} jatuh ketika \textit{node} 6 mati, sebab
\textit{node} itu kebetulan dititipi tujuh dari 26 kunci. Ketimpangan titipan tersebut
adalah kelemahan Chord yang sudah dikenal, dan Chord menjawabnya dengan \textit{node}
maya beserta penyalinan kunci ke beberapa penerus~\cite{stoica2001chord}.

\section{Pertumbuhan Saat Jaringan Membesar}

Delapan \textit{node} terlalu sedikit untuk memutuskan apa pun. Pertanyaan yang
menentukan adalah apa yang terjadi ketika jumlah \textit{node} digandakan. Bila ongkos
pesan tumbuh sebanding dengan jumlah \textit{node}, pertumbuhannya 2,00 kali. Bila
tumbuh sebanding dengan logaritmanya, pertumbuhannya 1,33 kali.

Tabel~\ref{tab:penskalaan-p2p} memuat hasilnya. Angkanya diambil dari keluaran
\texttt{ukur\_penskalaan.py}.

\begin{table}[htbp]
  \centering
  \caption{Pertumbuhan pesan ketika jumlah \textit{node} dinaikkan dari 8 ke 16}
  \label{tab:penskalaan-p2p}
  \begin{tabularx}{\textwidth}{|X|l|l|l|}
    \hline
    \textbf{Yang diukur} & \textbf{Client-server} & \textbf{\textit{Unstructured}} & \textbf{\textit{Structured}} \\
    \hline
    Pesan median 8 \textit{node} & 2 & 19 & 4 \\
    \hline
    Pesan median 16 \textit{node} & 2 & 35 & 6 \\
    \hline
    Pertumbuhan pesan & 1,00 kali & 1,84 kali & 1,50 kali \\
    \hline
    \textit{Hop} median 8 \textit{node} & 1 & 2 & 2 \\
    \hline
    \textit{Hop} median 16 \textit{node} & 1 & 2 & 3 \\
    \hline
    Terjawab pada 16 \textit{node} & 200 & 199 & 200 \\
    \hline
  \end{tabularx}
\end{table}

Pertumbuhan susunan \textit{unstructured}, yaitu 1,84 kali, sangat dekat dengan
pertumbuhan jumlah \textit{node}-nya. Pertumbuhan susunan \textit{structured}, yaitu 1,50 kali,
berada di antara 1,33 dan 2,00. Angka 1,50 itu bukan penyimpangan dari teori,
melainkan akibat \textit{hop} yang hanya dapat bernilai bulat. \textit{Hop} mediannya
naik dari dua menjadi tiga, dan pesannya memang dua kali \textit{hop}, sehingga
empat menjadi enam.

Baris terakhir memuat kejanggalan yang penghabisan. Satu \textit{lookup} pada susunan
\textit{unstructured} tidak terjawab pada jaringan enam belas \textit{node}, padahal seluruh
\textit{node}-nya hidup. Sebabnya batas empat \textit{hop} tidak lagi menjangkau seluruh
jaringan. Menaikkan batas itu memang memulihkan jawabannya, tetapi sekaligus
menaikkan jumlah pesannya, dan pertukaran itulah yang tidak dimiliki susunan
\textit{structured}.

\section{Konteks Pemakaian Ketiga Susunan}
\label{sec:konteks-p2p}

Ketiga susunan bukan selera, melainkan jawaban atas keadaan yang berbeda. Dua
pertanyaan menentukan pilihannya. Pertanyaan pertama, apakah ada pihak yang
berwenang memegang indeks. Pertanyaan kedua, seberapa sering keanggotaan
jaringannya berubah.

\textit{Centralized index} dipakai ketika satu pihak memang memiliki dan
mengawasi jaringannya, keanggotaannya jarang berubah, dan ongkos \textit{lookup}
harus murah sekaligus dapat diramalkan. Napster memakainya untuk daftar
\textit{file}, dan \textit{tracker} BitTorrent memakai bentuk yang sama sampai
hari ini~\cite{cohen2008bittorrent}. Tanpa indeks itu, setiap \textit{lookup}
harus dicari, yaitu 19 pesan pada contoh bab ini, bukan dua. Harganya satu, yaitu
titik kegagalan tunggal, sebab nol persen \textit{lookup} terjawab ketika
indeksnya mati.

\textit{Query flooding} dipakai ketika tidak ada pihak yang berwenang menyimpan
indeks, dan \textit{node} bebas datang serta pergi. Gnutella lahir dari keadaan
itu sesudah Napster dihentikan pengadilan, dan jaringannya bekerja tanpa satu
pun server pusat~\cite{ripeanu2001gnutella}. Tanpa cara ini, jaringan tanpa
pemilik tidak punya jalan menemukan apa pun kecuali mengangkat sebuah indeks,
dan indeks itu kembali menjadi titik kegagalan tunggal. Harganya ongkos pesan
yang tumbuh hampir sebanding dengan jumlah \textit{node}, yaitu 1,84 kali ketika
\textit{node}-nya digandakan, ditambah \textit{lookup} yang tidak terjawab
ketika batas \textit{hop} tidak lagi menjangkau seluruh jaringan.

\textit{Identifier ring} dipakai ketika jaringannya besar, keanggotaannya
berubah terus, dan \textit{lookup} harus tetap murah sekaligus selalu terjawab.
Amazon Dynamo menempatkan data keranjang belanja dengan \textit{consistent
hashing} di atas \textit{ring}, sebab mesin di pusat datanya datang dan pergi
setiap hari~\cite{decandia2007dynamo}. Tanpa cara ini, pilihan yang tersisa
hanya indeks terpusat atau penyebaran, dan keduanya sudah punya harga
masing-masing. Harganya tabel tambahan di setiap \textit{node}, yaitu
\textit{finger table}, beserta pemeliharaannya ketika \textit{node} berpindah.

Susunan \textit{structured} tidak dapat menjawab pertanyaan seperti itu lewat
satu \textit{lookup}. Sebabnya letak kunci dihitung dari hash nama yang lengkap,
sedangkan hash potongan nama menunjuk titik yang tidak ada hubungannya dengan
kunci mana pun. Tabel~\ref{tab:sebagian} memuat hasil percobaannya pada potongan
nama USD, dan angkanya diambil dari keluaran
\texttt{uji\_pencarian\_sebagian.py}.

\begin{table}[htbp]
  \centering
  \caption{Pertanyaan yang hanya menyebut sebagian nama, yaitu USD}
  \label{tab:sebagian}
  \begin{tabularx}{\textwidth}{|X|l|l|}
    \hline
    \textbf{Susunan} & \textbf{Kunci ketemu} & \textbf{Pesan} \\
    \hline
    Client-server & 1 dari 1 & 2 \\
    \hline
    \textit{Unstructured} & 1 dari 1 & 14 \\
    \hline
    \textit{Structured} & 0 dari 1 & 2 \\
    \hline
  \end{tabularx}
\end{table}

Baris terakhir tabel tersebut perlu dibaca dengan tepat. Susunan
\textit{structured} bukan hanya gagal menemukan, melainkan memang tidak dirancang
untuk pertanyaan itu, sebab yang dijawabnya hanya kunci yang disebut selengkapnya.
Dua pesan yang telanjur dipakainya karena itu terbuang. Pilihan yang tersisa bagi
susunan \textit{structured} adalah bertanya kepada seluruh \textit{node}, dan
pada saat itu ongkosnya sama saja dengan penyebaran.

Tabel~\ref{tab:konteks-p2p} merangkum ketiganya, dan angkanya diambil dari
pengukuran bab ini.

\begin{table}[htbp]
  \centering
  \caption{Keadaan yang menentukan pilihan susunan beserta harganya}
  \label{tab:konteks-p2p}
  \begin{tabularx}{\textwidth}{|l|X|X|}
    \hline
    \textbf{Susunan} & \textbf{Dipakai ketika} & \textbf{Harga yang dibayar} \\
    \hline
    Client-server & Jaringan dimiliki satu pihak, keanggotaan jarang berubah, ongkos harus dapat diramalkan & Nol persen \textit{lookup} terjawab ketika indeksnya mati \\
    \hline
    \textit{Unstructured} & Tidak ada pihak yang berwenang memegang indeks, \textit{node} bebas datang dan pergi & 19 pesan per \textit{lookup}, dan tumbuh 1,84 kali ketika \textit{node} digandakan \\
    \hline
    \textit{Structured} & Jaringan besar, keanggotaan berubah terus, \textit{lookup} harus murah dan selalu terjawab & \textit{Finger table} di setiap \textit{node} beserta pemeliharaannya \\
    \hline
  \end{tabularx}
\end{table}

Satu catatan praktis menutup seksi ini. Ketiga susunan tersebut dapat hidup
berdampingan di dalam satu sistem. BitTorrent memakai \textit{tracker} sebagai
indeks terpusat, sekaligus menyediakan \textit{distributed hash table}
berdasarkan Kademlia bagi klien yang \textit{tracker}-nya tidak dapat
dihubungi~\cite{maymounkov2002kademlia}. Pilihannya karena itu bukan satu
susunan untuk selamanya, melainkan susunan mana untuk bagian mana.

\section{Ringkasan}

\textit{Peer-to-peer} menata sistem menjadi \textit{node} yang berkedudukan sama,
sehingga tidak ada pihak yang selalu melayani. Kedudukan yang sama itu membawa
satu pekerjaan baru, yaitu menemukan data tanpa ada satu pun \textit{node} yang
mengetahui seluruh isi jaringan. Seluruh perbedaan yang diukur bab ini
berpangkal pada pekerjaan tersebut.

Ketiga susunan menjawabnya dengan pengetahuan yang berbeda. \textit{Centralized index}
menyimpan seluruh pengetahuan pada satu tempat, lalu membayarnya dengan dua
pesan per \textit{lookup}. Susunan \textit{unstructured} tidak menyimpannya di mana pun,
lalu membayarnya dengan 19 pesan. Susunan \textit{structured} menghitung letaknya dari
kuncinya sendiri, dan cukup membayar empat pesan.

Harga sebenarnya baru terlihat ketika satu \textit{node} mati. Rata-rata ketiga
susunan hampir sama, yaitu berkisar 83 sampai 88 persen, sehingga rata-rata itu
tidak berguna. Yang membedakan adalah giliran terburuknya, yaitu nol persen
pada client-server melawan 55,2 dan 71,5 persen pada kedua susunan lain.
Pelajaran itu sama dengan Bab~2, yaitu ujung sebaran lebih banyak bercerita
daripada nilai tengahnya.

Penggandaan jumlah \textit{node} memisahkan kedua susunan \textit{peer-to-peer}. Pesan
susunan \textit{unstructured} tumbuh 1,84 kali, hampir sebanding dengan jumlah
\textit{node}-nya. Pesan susunan \textit{structured} tumbuh 1,50 kali, dan \textit{hop}-nya hanya
bertambah satu. Selisih itu masih kecil pada enam belas \textit{node}, sedangkan pada
jutaan \textit{node} selisih yang sama menjadi jurang.

Pesan utama bab ini karena itu bukan bahwa \textit{peer-to-peer} mengungguli
client-server. Pesannya adalah bahwa ketiadaan pusat harus dibayar dengan
pekerjaan \textit{lookup}, dan besar pembayarannya ditentukan oleh seberapa banyak
struktur yang bersedia dipelihara. Tanpa struktur, pembayarannya disebar ke
seluruh \textit{neighbor}. Dengan struktur, pembayarannya dihitung.

\section{Latihan}

Ketiga latihan berikut memakai satu basis kode yang sama di
\berkas{sources/chapter07/}, tetapi menyasar tiga masalah yang berbeda.
Latihan~1 menghitung \textit{lookup cost}. Latihan~2 mengukur
\textit{failure resilience} ketika satu \textit{node} mati. Latihan~3 memeriksa
\textit{message scalability} ketika jaringannya membesar. Hasil setiap mahasiswa dapat berbeda bila \textit{seed}-nya diubah, dan yang
dilaporkan adalah pola perbandingannya, bukan angka mutlaknya.

Venv yang sama dipakai kembali, dan pengaktifannya seperti pada
Listing~\ref{lst:venv-ch07}. Bab ini hanya memakai pustaka standar Python,
sehingga tidak ada pemasangan tambahan.

\begin{lstlisting}[language=bash, caption={Mengaktifkan venv sebelum ketiga
  latihan, skrip tersedia di \berkas{sources/siapkan_venv.sh}},
  label=lst:venv-ch07]
bash sources/siapkan_venv.sh
source .venv/bin/activate
cd sources/chapter07
python susunan.py daftar
# Keluarannya baris pertama: Kunci kurs terdaftar : 26 (18 September 2026)
python susunan.py cari cincin EUR/IDR
# Keluarannya baris terakhir: Jumlah lompatan : 3
\end{lstlisting}

\subsection*{Latihan 1: \textit{Lookup Cost}}

Latihan ini membandingkan ongkos satu \textit{lookup} pada susunan \textit{unstructured}
dengan susunan \textit{structured}. Keduanya memakai jumlah sambungan yang sama, yaitu
dua belas, sehingga selisihnya benar-benar berasal dari cara mencari. Langkah
kerjanya sebagai berikut.

\begin{enumerate}
  \item Jalankan ketiga perintah pada Listing~\ref{lst:jalankan-ukur-pencarian},
    lalu catat keenam angka sebaran untuk setiap susunan.
  \item Buka \berkas{sources/chapter07/tetangga.py}, lalu tunjukkan baris yang
    membuang pertanyaan berulang.
  \item Ubah \texttt{BATAS\_LOMPATAN} pada \textit{file} itu dari empat menjadi dua,
    jalankan ulang, lalu catat jumlah pesan beserta \textit{lookup} yang terjawab.
    Kembalikan angkanya sesudah dicatat.
  \item Jalankan \texttt{python susunan.py cari tetangga EUR/IDR}, lalu
    bandingkan angkanya dengan \texttt{python susunan.py cari cincin EUR/IDR}.
  \item Jalankan \texttt{python uji\_pencarian\_sebagian.py}, lalu catat berapa
    kunci yang ketemu pada ketiga susunan beserta pesannya.
  \item Buka \berkas{sources/chapter07/cincin.py}, ganti sementara baris
    \texttt{berikut = self.jari\_terdekat(sekarang, sasaran) or penerus}
    menjadi \texttt{berikut = penerus}, jalankan ulang pengukurannya, lalu catat
    \textit{hop} dan pesannya. Penggantian itu mematikan \textit{finger table},
    sehingga penelusurannya merayap lewat penerus saja. Kembalikan barisnya
    sesudah dicatat.
\end{enumerate}

\begin{lstlisting}[language=bash, caption={Mengukur \textit{lookup cost},
  skrip tersedia di \berkas{sources/chapter07/ukur_pencarian.py}},
  label=lst:jalankan-ukur-pencarian]
python ukur_pencarian.py
# Keluarannya pada susunan unstructured: Pesan median     : 19.0
#                                        Hop median       : 2.0
#                                        Lookup terjawab  : 200 dari 200
python susunan.py cari tetangga EUR/IDR
# Keluarannya: Jumlah pesan    : 20
#              Jumlah hop      : 3
python susunan.py cari cincin EUR/IDR
# Keluarannya: Jumlah pesan    : 6
#              Jumlah hop      : 3
# Sesudah finger table dimatikan: Jumlah pesan : 12
#                                 Jumlah hop   : 6
python uji_pencarian_sebagian.py
# Keluarannya: Client-server ketemu       : 1 (2 pesan)
#              Unstructured ketemu        : 1 (14 pesan)
#              Structured ketemu          : 0 (2 pesan)
\end{lstlisting}

Tabel~\ref{tab:lat1p2p-contoh} berisi hasil satu kali pengukuran nyata.

\begin{table}[htbp]
  \centering
  \caption{Contoh hasil Latihan 1 yang sudah terisi, dua ratus \textit{lookup}}
  \label{tab:lat1p2p-contoh}
  \begin{tabularx}{\textwidth}{|X|l|l|l|}
    \hline
    \textbf{Yang dilaporkan} & \textbf{Client-server} & \textbf{\textit{Unstructured}} & \textbf{\textit{Structured}} \\
    \hline
    Pesan median & 2 & 19 & 4 \\
    \hline
    Pesan persentil ke-95 & 2 & 20 & 6 \\
    \hline
    Pesan terbesar & 2 & 20 & 8 \\
    \hline
    \textit{Hop} median & 1 & 2 & 2 \\
    \hline
    \textit{Hop} terbesar & 1 & 3 & 4 \\
    \hline
    Sisi jaringan & 7 & 12 & 13 \\
    \hline
    Pesan median dengan batas dua & 2 & 9 & 4 \\
    \hline
    Terjawab dengan batas dua & 200 & 176 & 200 \\
    \hline
    \textit{Hop} median tanpa \textit{finger table} & 1 & 2 & 4 \\
    \hline
  \end{tabularx}
\end{table}

Dua baris terakhir tabel tersebut memuat pertukaran yang menjadi inti latihan
ini. Menurunkan batas \textit{hop} menjadi dua memangkas pesan susunan
\textit{unstructured} dari 19 menjadi 9, tetapi 24 \textit{lookup} lalu tidak
terjawab.
Susunan \textit{structured} tidak terpengaruh sama sekali, sebab penelusurannya tidak
pernah bergantung pada batas itu.

Isikan hasil pengukuran pada Tabel~\ref{tab:lat1p2p-hasil}.

\begin{table}[htbp]
  \centering
  \caption{Lembar isian Latihan 1}
  \label{tab:lat1p2p-hasil}
  \begin{tabularx}{\textwidth}{|X|l|l|l|}
    \hline
    \textbf{Yang dilaporkan} & \textbf{Client-server} & \textbf{\textit{Unstructured}} & \textbf{\textit{Structured}} \\
    \hline
    Pesan median & \dots & \dots & \dots \\
    \hline
    Pesan terbesar & \dots & \dots & \dots \\
    \hline
    \textit{Hop} terbesar & \dots & \dots & \dots \\
    \hline
    Pesan dengan batas dua & \dots & \dots & \dots \\
    \hline
    Terjawab dengan batas dua & \dots & \dots & \dots \\
    \hline
    \textit{Hop} median tanpa \textit{finger table} & \dots & \dots & \dots \\
    \hline
  \end{tabularx}
\end{table}

Jawab tiga pertanyaan berikut dengan menunjuk angka pada tabel. Pertama,
jelaskan mengapa sebaran pesan susunan \textit{unstructured} hampir rata, sedangkan
sebaran susunan \textit{structured} melebar dari empat sampai delapan. Kedua, sebutkan
apa yang diperoleh dan apa yang hilang ketika batas \textit{hop} diturunkan menjadi
dua. Ketiga, bandingkan jumlah \textit{hop} susunan \textit{structured} dengan dan
tanpa \textit{finger table}, lalu jelaskan mengapa jalan pintas itu dibutuhkan
padahal letak kuncinya sudah dapat dihitung.

\subsection*{Latihan 2: \textit{Failure Resilience}}

Latihan ini mengukur berapa \textit{lookup} yang masih terjawab ketika satu \textit{node}
keluar dari jaringan. Di sinilah perbedaan client-server dan
\textit{peer-to-peer} muncul sebagai angka, bukan sebagai pendapat. Langkah
kerjanya sebagai berikut.

\begin{enumerate}
  \item Jalankan ketiga perintah pada Listing~\ref{lst:jalankan-uji-ketahanan},
    lalu catat kelima baris hasil untuk setiap susunan.
  \item Bandingkan baris rata-rata dengan baris terburuk pada susunan
    client-server, lalu catat selisihnya.
  \item Jalankan \texttt{python tetangga.py}, lalu cari \textit{node} yang hanya punya
    satu \textit{neighbor} beserta jumlah kunci yang dipegangnya.
  \item Jalankan \texttt{python cincin.py}, lalu cari \textit{node} yang dititipi
    kunci paling banyak.
\end{enumerate}

\begin{lstlisting}[language=bash, caption={Mengukur \textit{failure
  resilience}, skrip tersedia di
  \berkas{sources/chapter07/uji_ketahanan.py}}, label=lst:jalankan-uji-ketahanan]
python uji_ketahanan.py
# Keluarannya pada susunan client-server: Node 0 mati             : 0.0 persen
#                                         Terburuk satu node mati : 0.0 persen
python tetangga.py
# Keluarannya memuat: node 5 neighbor [1] kunci 6
python cincin.py
# Keluarannya memuat: node 6 pengenal 27881 jari [...] kunci 7
\end{lstlisting}

Tabel~\ref{tab:lat2p2p-contoh} berisi hasil satu kali percobaan nyata.

\begin{table}[htbp]
  \centering
  \caption{Contoh hasil Latihan 2 yang sudah terisi, dua ratus \textit{lookup}}
  \label{tab:lat2p2p-contoh}
  \begin{tabularx}{\textwidth}{|X|l|l|l|}
    \hline
    \textbf{Yang dilaporkan} & \textbf{Client-server} & \textbf{\textit{Unstructured}} & \textbf{\textit{Structured}} \\
    \hline
    \textit{Node} 0 mati & 0,0 persen & 94,0 persen & 86,0 persen \\
    \hline
    Rata-rata satu \textit{node} mati & 87,5 persen & 83,5 persen & 87,9 persen \\
    \hline
    Terburuk satu \textit{node} mati & 0,0 persen & 55,2 persen & 71,5 persen \\
    \hline
    Nomor \textit{node} terburuk & 0 & 1 & 6 \\
    \hline
    Kunci pada \textit{node} terburuk & 26 & 4 & 7 \\
    \hline
  \end{tabularx}
\end{table}

Isikan hasil percobaan pada Tabel~\ref{tab:lat2p2p-hasil}.

\begin{table}[htbp]
  \centering
  \caption{Lembar isian Latihan 2}
  \label{tab:lat2p2p-hasil}
  \begin{tabularx}{\textwidth}{|X|l|l|l|}
    \hline
    \textbf{Yang dilaporkan} & \textbf{Client-server} & \textbf{\textit{Unstructured}} & \textbf{\textit{Structured}} \\
    \hline
    \textit{Node} 0 mati & \dots & \dots & \dots \\
    \hline
    Rata-rata satu \textit{node} mati & \dots & \dots & \dots \\
    \hline
    Terburuk satu \textit{node} mati & \dots & \dots & \dots \\
    \hline
    Nomor \textit{node} terburuk & \dots & \dots & \dots \\
    \hline
  \end{tabularx}
\end{table}

Jawab tiga pertanyaan berikut. Pertama, rata-rata susunan client-server justru
paling tinggi pada tabel tersebut. Jelaskan mengapa angka itu menyesatkan, lalu
sebutkan baris mana yang seharusnya dibaca. Kedua, jelaskan mengapa matinya
\textit{node} 1 pada susunan \textit{unstructured} lebih merugikan daripada matinya \textit{node}
lain, padahal kunci yang dipegangnya hanya empat. Ketiga, sebutkan satu
perubahan yang akan menaikkan angka terburuk pada susunan \textit{structured}, beserta
ongkos yang harus dibayar untuk perubahan itu.

\subsection*{Latihan 3: \textit{Message Scalability}}

Latihan ini memeriksa apakah pertumbuhan pesan sebanding dengan jumlah \textit{node}
atau dengan logaritmanya. Jumlah \textit{node} dinaikkan dari delapan menjadi enam
belas, sehingga pertumbuhan yang sebanding dengan jumlah \textit{node} bernilai 2,00
kali. Langkah kerjanya sebagai berikut.

\begin{enumerate}
  \item Jalankan kedua perintah pada Listing~\ref{lst:jalankan-penskalaan}, lalu
    catat kedua angka pembanding pada kepala laporannya.
  \item Catat pertumbuhan pesan ketiga susunan beserta \textit{hop} mediannya.
  \item Jalankan \texttt{python ukur\_pencarian.py 32}, lalu catat pesan
    mediannya pada ketiga susunan.
  \item Catat berapa \textit{lookup} yang tidak terjawab pada susunan
    \textit{unstructured}, lalu jelaskan sebabnya.
\end{enumerate}

\begin{lstlisting}[language=bash, caption={Mengukur \textit{message
  scalability}, skrip tersedia di
  \berkas{sources/chapter07/ukur_penskalaan.py}}, label=lst:jalankan-penskalaan]
python ukur_penskalaan.py
# Keluarannya: Pertumbuhan jumlah node    : 2.00 kali
#              Pertumbuhan logaritmanya   : 1.33 kali
# Keluarannya pada susunan unstructured: Pertumbuhan pesan : 1.84 kali
# Keluarannya pada susunan structured  : Pertumbuhan pesan : 1.50 kali
python ukur_pencarian.py 32
# Keluarannya pada susunan unstructured: Pesan median    : 52.0
#                                        Lookup terjawab : 172 dari 200
\end{lstlisting}

Tabel~\ref{tab:lat3p2p-contoh} berisi hasil satu kali pengukuran nyata.

\begin{table}[htbp]
  \centering
  \caption{Contoh hasil Latihan 3 yang sudah terisi, dua ratus \textit{lookup} per
    ukuran}
  \label{tab:lat3p2p-contoh}
  \begin{tabularx}{\textwidth}{|X|l|l|l|}
    \hline
    \textbf{Yang dilaporkan} & \textbf{Client-server} & \textbf{\textit{Unstructured}} & \textbf{\textit{Structured}} \\
    \hline
    Pesan median 8 \textit{node} & 2 & 19 & 4 \\
    \hline
    Pesan median 16 \textit{node} & 2 & 35 & 6 \\
    \hline
    Pertumbuhan pesan & 1,00 kali & 1,84 kali & 1,50 kali \\
    \hline
    \textit{Hop} median 16 \textit{node} & 1 & 2 & 3 \\
    \hline
    Terjawab pada 16 \textit{node} & 200 & 199 & 200 \\
    \hline
    Pesan median 32 \textit{node} & 2 & 52 & 6 \\
    \hline
    Terjawab pada 32 \textit{node} & 200 & 172 & 200 \\
    \hline
  \end{tabularx}
\end{table}

Pembandingnya dua angka, yaitu 2,00 kali bila ongkos tumbuh sebanding dengan
jumlah \textit{node}, dan 1,33 kali bila tumbuh sebanding dengan logaritmanya. Isikan
hasil pengukuran pada Tabel~\ref{tab:lat3p2p-hasil}.

\begin{table}[htbp]
  \centering
  \caption{Lembar isian Latihan 3}
  \label{tab:lat3p2p-hasil}
  \begin{tabularx}{\textwidth}{|X|l|l|l|}
    \hline
    \textbf{Yang dilaporkan} & \textbf{Client-server} & \textbf{\textit{Unstructured}} & \textbf{\textit{Structured}} \\
    \hline
    Pesan median 8 \textit{node} & \dots & \dots & \dots \\
    \hline
    Pesan median 16 \textit{node} & \dots & \dots & \dots \\
    \hline
    Pesan median 32 \textit{node} & \dots & \dots & \dots \\
    \hline
    Pertumbuhan pesan & \dots & \dots & \dots \\
    \hline
    Terjawab pada 16 \textit{node} & \dots & \dots & \dots \\
    \hline
  \end{tabularx}
\end{table}

Jawab tiga pertanyaan berikut dengan menunjuk angka pada tabel. Pertama,
tentukan susunan mana yang pertumbuhannya paling dekat dengan 2,00 kali, lalu
jelaskan sebabnya dari cara kerjanya. Kedua, jelaskan mengapa pertumbuhan
susunan \textit{structured} tercatat 1,50 kali, bukan 1,33 kali seperti yang
diperkirakan. Ketiga, dengan menggabungkan hasil ketiga latihan, tentukan
susunan mana yang akan dipakai untuk jaringan berisi sejuta \textit{node}, lalu
sebutkan tiga angka yang menjadi dasar keputusannya.
